\section{Intersection Point of Two Linked Lists}
\label{sec:ll-intersection}

\problemstatement{Given the heads of two singly linked lists that may merge
at some node (forming a Y shape), return the first common node, or null if
they never intersect. Use $O(1)$ extra space.}

\begin{patternbox}
The elegant solution is the \textbf{two-pointer switch}: walk one pointer
along list A then list B, and another along list B then list A. They travel
the same total distance, so they arrive at the intersection simultaneously
(or at null together).
\end{patternbox}

\subsection*{Equalise the path lengths automatically}

Let the lists have lengths $a$ and $b$ with a shared tail of length $c$.
Pointer 1 walks $a$ then $b$; pointer 2 walks $b$ then $a$; both cover
$a+b$ nodes. After the switch, each is exactly $c$ nodes from the end, so
they meet at the first shared node. If there is no intersection, both reach
null at the same time. This removes the need to measure and align lengths
manually.

\begin{ideabox}[Switching Lists Cancels the Length Difference]
By having each pointer traverse both lists, the head-start difference $|a-b|$
is absorbed: after the switch they are equidistant from the merge point. The
first place they coincide is the intersection.
\end{ideabox}

\subsection*{Algorithm}

\begin{enumerate}
  \item \texttt{p = headA}, \texttt{q = headB}.
  \item Advance both; when \texttt{p} hits null, redirect it to
        \texttt{headB}; when \texttt{q} hits null, redirect to
        \texttt{headA}.
  \item Stop when \texttt{p == q}; that node (or null) is the answer.
\end{enumerate}

\subsection*{Dry run}

\dryrun{A $=1\to9\to7$, B $=6\to7$, sharing the node $7$.}

\begin{itemize}
  \item p: $1,9,7,\text{(B)}6,7$. q: $6,7,\text{(A)}1,9,7$.
  \item Both reach the shared $7$ at the same step -- the intersection.
\end{itemize}

\subsection*{C++ implementation}

\begin{lstlisting}
#include <bits/stdc++.h>
using namespace std;

struct ListNode {
    int val; ListNode* next;
    ListNode(int v) : val(v), next(nullptr) {}
};

ListNode* getIntersectionNode(ListNode* a, ListNode* b) {
    if (!a || !b) return nullptr;
    ListNode* p = a; ListNode* q = b;
    while (p != q) {
        p = p ? p->next : b;                      // switch to B at the end
        q = q ? q->next : a;                      // switch to A at the end
    }
    return p;                                     // meeting node or null
}

int main() {
    ListNode* shared = new ListNode(7);
    shared->next = new ListNode(8);
    ListNode* a = new ListNode(1);
    a->next = new ListNode(9); a->next->next = shared;
    ListNode* b = new ListNode(6); b->next = shared;

    ListNode* r = getIntersectionNode(a, b);
    cout << (r ? r->val : -1) << "\n";             // 7
    return 0;
}
\end{lstlisting}

\begin{complexitybox}
\textbf{Time Complexity.} $O(a+b)$ -- each pointer walks both lists once.
\textbf{Space Complexity.} $O(1)$ (a hash set would cost $O(n)$).
\end{complexitybox}

\begin{takeawaybox}
The two-pointer list switch equalises unequal lengths automatically and
finds the intersection in constant space. It is a favourite interview
trick precisely because the length-cancellation is so clean.
\end{takeawaybox}
